\chapter{The frame scheduler}
\label{ch:scheduler}

\begin{aside}
A TUI is event-driven, not frame-driven. There is no 60 Hz
refresh; there is no v-sync. Frames happen because something
asked for one. This chapter is about who asks, and how often.
\end{aside}

\section{What is a frame?}

A frame in \maya{} is a single render pass: clear back buffer,
paint, diff, emit. Frames are not free --- each is at least one
\code{write} syscall and a few microseconds of CPU --- so we want
exactly as many as the application needs and no more.

\section{Triggers}

Three things trigger a frame:

\begin{enumerate}[leftmargin=*]
  \item \textbf{Input.} A key press, a mouse event, a paste, a
    focus change, a resize.
  \item \textbf{Signal change.} Application state updated via
    the signal graph (see Chapter~\ref{ch:signals}).
  \item \textbf{Animation.} A widget registered an ongoing
    animation that wants a frame at a specific time.
\end{enumerate}

A frame triggered by any source serves all of them: if input and
a signal change arrive together, one frame paints both.

\section{The event loop}

The loop is roughly:

\begin{lstlisting}
while (running) {
  Event ev = wait_for_next_event(deadline);
  dispatch(ev);
  if (frame_dirty) {
    render_frame();
    frame_dirty = false;
  }
}
\end{lstlisting}

\code{wait\_for\_next\_event} sleeps the thread until input
arrives or the deadline fires. The deadline is the soonest of:

\begin{itemize}[leftmargin=*]
  \item Animation frame deadline (e.g. ``redraw at $t + 16$ ms''
    for a 60-Hz animation).
  \item Timeout (a registered timer).
  \item Forever, if neither of the above.
\end{itemize}

\section{Coalescing}

Multiple events in a single wakeup are dispatched in order
before rendering. This coalesces:

\begin{itemize}[leftmargin=*]
  \item A paste of 500 bytes: 500 keys, one frame.
  \item A burst of mouse-move events: many moves, one frame.
  \item A flurry of signal changes during a command: many
    updates, one frame.
\end{itemize}

Coalescing makes the renderer's throughput O(unique frames), not
O(events).

\section{Frame budget}

Sixty frames per second gives 16.7 ms per frame. \maya{}'s
target is to finish a frame in under 5 ms, leaving headroom for
the application's own work.

Of that 5 ms:

\begin{itemize}[leftmargin=*]
  \item Layout: 0.5--1 ms typical, depends on tree depth.
  \item Paint: 0.5--2 ms, scales with cell count.
  \item Diff and emit: 0.5--1 ms.
  \item Application: whatever is left.
\end{itemize}

If a frame goes over budget, the next event is processed
immediately and a new frame is requested. There is no
back-pressure: the renderer cannot ask the keyboard to slow down.

\section{Animation tick}

An animation registers a callback and a duration:

\begin{lstlisting}
animate(300ms, [](float t) {
    // t goes from 0.0 to 1.0
    ctx.set(progress, t);
});
\end{lstlisting}

The scheduler picks a tick rate (default 60 Hz) and fires the
callback every tick with a $t$ value. When $t$ reaches 1.0 the
animation completes.

Many animations can be active at once; the scheduler picks the
deadline of the soonest tick and wakes then.

\section{Tick deduplication}

If three animations all want a tick at $t + 16$ ms, the
scheduler wakes once and fires all three. One frame paints the
combined effect.

\section{Debouncing}

For events that fire faster than the scheduler can keep up
(rapid scrolling, mouse motion), \maya{} debounces by simply
processing the latest. The wakeup empties the input queue and
dispatches all events; the renderer paints once with the final
state.

\begin{tryit}
Hold a key down. The repeat rate of your OS may be 30 Hz or
faster, but \maya{} processes it as ``keys'' not as ``frames''
--- one frame per repeat is fine, but a slow widget might
combine several into one.
\end{tryit}

\section{Sleeping}

When idle, \maya{} blocks on \code{poll} (Linux),
\code{kqueue} (BSD), or \code{select}. The thread sleeps; CPU
goes to 0\%. There is no spinning, no busy wait.

\begin{funfact}
A long-running TUI like a chat client uses about 0.1\% CPU when
idle. Compare to a browser tab with the same chat, which often
uses 5--10\% for the cursor blink alone.
\end{funfact}

\section{Resize handling}

SIGWINCH (window change signal) interrupts the syscall and
returns \code{EINTR}. \maya{}'s loop catches this, queries the
new terminal size via \code{TIOCGWINSZ}, and forces a full
redraw.

The redraw is forced via the sentinel trick from
Chapter~\ref{ch:buffers}: the front buffer is filled with sentinels,
the next frame diffs every cell.

\section{Suspend and resume}

Ctrl-Z sends SIGTSTP. \maya{} catches it, restores the terminal
to cooked mode, prints the appropriate ANSI to exit alt-screen,
and re-raises SIGTSTP so the shell can suspend the process.

On SIGCONT the loop reverses the dance: alt-screen, raw mode,
sentinel-flood, render.

\begin{warning}
If a custom shutdown handler raises an exception, the terminal
may be left in raw mode. \maya{} installs handlers in the
constructor with RAII guards so even \code{std::terminate}
leaves the terminal sane.
\end{warning}

\section{The render-on-every-tick fallacy}

Some libraries render at a fixed rate regardless of state:
``redraw at 60 Hz''. This is wasteful for a TUI. A file
browser sitting on the same folder for ten minutes should
emit zero ANSI in those ten minutes.

\maya{}'s scheduler emits ANSI only when something changed. If
nothing did, the screen is byte-for-byte the same as before.

\section{Tearing}

Without DEC 2026 synchronised output, a long frame may be
emitted in chunks. The user can see a half-painted screen for
one tick. \maya{} mitigates by sending the smallest possible
runs and by using \code{TCSADRAIN} on resize.

DEC 2026, when supported, brackets the whole frame so the
terminal applies it atomically.

\section{Recap}

\begin{recap}
\begin{itemize}[leftmargin=*]
  \item Frames are demand-driven: input, signals, or
    animation.
  \item The loop coalesces events: many events, one frame.
  \item Idle costs zero CPU.
  \item Animation tick rate defaults to 60 Hz, configurable.
  \item Resize and suspend handlers restore the terminal.
\end{itemize}
\end{recap}

\section*{Workshop}

\begin{enumerate}[leftmargin=*]
  \item Add a frame counter to your status bar. Confirm it
    increments only when something changes.
  \item Register an animation that fades a colour over 1
    second. Observe the tick rate via the counter.
  \item Suspend and resume your app via Ctrl-Z and \code{fg}.
    The terminal should restore cleanly each time.
\end{enumerate}
